\subsection{The \(B_{19}\) long head and its paid gains}
\label{app:exact19:Bhead}

The actual long head is isolated.
Indeed deleting one retains
\(B_{18}\), whose complete finite
one-node list is isolation or
simple-three at its long leaf two.
A second contact makes a cycle;
another four or an additional fork
is not finite; exceptional finite
ranks are too small.
The nonisolated cap would have
\[
 v=\gamma\omega_1+\omega_2/2,\quad
 A^2=2\gamma^2+2\gamma,\quad\cosh L_1=2\gamma+1 .
\]
Own zero has zero time and own
three is the actual affine
triangle, so \(\gamma\ge1/\sqrt2\)
and \(x_1>\sqrt{\sqrt2-1}>.643\).

The OLD actual bulk through fifteen
has some strong \(K=14\) or fifteen.
If old five is strong, its paid
reference has energy \(>.90715\),
with every actual weak hole omitted.
Graft
\(.643\) into its first-five block
of mean \(<.361\).
The other four coordinates remain
\(>.2904\), above the outgoing
slope \(<.214\), and the gain is
\((5/4)(.643-.361)^2=.099405\).
If old five is weak, six is strong.
Begin only with the FORMAL
strong-five curve, whose scalar
energy is \(E_K^{\rm old}\).
Here \(f_i^{\rm old}=\sqrt i-.61/\sqrt2\)
and \(E_K^{\rm old}=(f_5^{\rm old})^2/5
+\sum_{i=5}^{K-1}d_i^2\).
Before assigning any pressure to
it, replace its ENTIRE first-six
block by \(.643\) and five equal
slopes of sum \(f_6^{\rm old}-.643\).
Their energy exceeds
\[
 .643^2+(2.01815-.643)^2/5=.7916565045 ,
\]
whereas the formal first-six
energy is \(E_6^{\rm old}<.6971\).
The new slopes are ordered:
\(.643>(f_6^{\rm old}-.643)/5>.275\),
above every outgoing slope
\(\le d_6<.197\).
There is NO coefficient at five.
All later actual weak knots lie
in \(7,\ldots,K-1\); at most four
are nonadjacent, each costing
\(<1/3200\).
Their affected open intervals
are disjoint from the first-six
replacement, even for a hole
at seven.
Only now combine the actual
pressures at one, six and the
remaining strong knots.
One square gives
\[
 \|x\|^2>
 E_K^{\rm old}-E_6^{\rm old}
      +.7916565045-4/3200
 >.9084-.6971+.7916565045-.00125
       =1.0017065045>1 .
\]
Both references are ordered and
paid by the original \(x\).
Thus the nonisolated head is
impossible.

For the isolated cap,
\(A^2=2\gamma^2-1>0\) excludes
own zero/three/four.
Own six has the affine \(6,3\)
principal on the actual first two
roots; parallel has its feasible
positive future null.
The remaining owns are exactly
native five, finite at least seven,
or dotted.
Their actual feet are
\eqref{app:exact19:eq:longfeet}.

A native-five head forces true
new strong five/six/seven.
For a helper at \(i\le7\),
\[
 c_L\cdot c_i=\gamma_0(2g+1)-A_0A_i,\qquad
       \gamma_0=\varphi/2,\quad A_0^2=(\varphi-1)/2 .
\]
Since \(\gamma_0^2/A_0^2=\varphi^3/2>21/10\),
\[
 A_i^2-\frac{\gamma_0^2}{A_0^2}(2g+1)^2
                   <(28g^2-32g-13)/5<0
                  \qquad(.8<g<1.2).
\]
The convex quadratic is negative
at both endpoints.
The actual mutual would be
positive.
For the remaining weak isolated
dotted five,
\[
 c_L\cdot c_5=\varphi g-A_0\sqrt{10g^2-1}.
\]
Its sign is strictly negative,
since the squared difference in
the opposite order is
\((4\varphi-6)g^2-(\varphi-1)/2
\ge(7\varphi-11)/2>0\).
Weakness gives \(g<1/.58<2\);
\(\varphi-\sqrt{5(\varphi-1)}>-3/20\)
then bounds the mutual in
\((-3/10,0)\), a forbidden native
interval.
This proves the stated three
strong sites without a later
terminal floor.

For any strong-five \(.58\)
reference with first-five mean
\(b=f_5/5<.3652\) and individual
\(P_5\) coefficient \(>.151\),
a head correction
\(\mu(Z-P_5/5)\) costs \(\mu/5\)
before its floor is used.
It is orthogonal to that constant
block and has square \(R-1/5\).
Its effective gain is
\[
                             G_L=2\mu(h-b)-\mu^2(R-1/5).
\]
The complete paid gains are
\[
\begin{array}{c|c|c}
 \text{head}&\mu&G_L>\\ \hline
 5&.56&.015456\\
 \text{finite }\ge7&.4&.02914\\
 \text{dotted}&\sqrt2/5&.03384 .
\end{array}
\tag{B19headgains}\label{app:exact19:eq:Bheadgains}
\]
The first uses \(h>.393\);
the second is
\eqref{app:exact19:eq:longF}
minus \(.8b\) plus \(.032\).
For the dotted row, expand its
gain as
\[
 .4\{\sqrt{2-u}-\sqrt{u(1-u)}-\sqrt2b\}
                                -.04(1.6-2u+u^2),
\]
where \(\sqrt2b<.5166\).
For \(u\le.5\), its lower bound is
\(.29504-.0916u-.4\sqrt{u(1-u)}>.04404\),
using \(\sqrt{.0916^2+.4^2}<.4104\).
For \(u\ge.5\), it is
\(.32912-.15976u-.4\sqrt{u(1-u)}>.03384\),
using \(\sqrt{.15976^2+.4^2}<.4308\).
The same two concave root chords
prove these inequalities.
All three individual prices fit.

